\section{Useful Things}
\subsection{Fast I/O}
\begin{lstlisting}
// C++
#pragma GCC optimize("O3")
ios_base::sync_with_stdio(0);
 cin.tie(0);

# Python
import sys
input = sys.stdin.readline
sys.stdout.write("------")
\end{lstlisting}
\subsection{Debug}
\begin{lstlisting}
void __print(int x) {cerr << x;}
void __print(long x) {cerr << x;}
void __print(long long x) {cerr << x;}
void __print(unsigned x) {cerr << x;}
void __print(unsigned long x) {cerr << x;}
void __print(unsigned long long x) {cerr << x;}
void __print(float x) {cerr << x;}
void __print(double x) {cerr << x;}
void __print(long double x) {cerr << x;}
void __print(char x) {cerr << '\'' << x << '\'';}
void __print(const char *x) {cerr << '\"' << x << '\"';}
void __print(const string &x) {cerr << '\"' << x << '\"';}
void __print(bool x) {cerr << (x ? "true" : "false");}
template<size_t N> void __print(const bitset<N> &x) {cerr << x;}
template<typename T, typename V> void __print(const pair<T,V> &x) {
    cerr << '{'; __print(x.first);
    cerr << ','; __print(x.second); cerr << '}';
}
template<typename T> void __print(const T &x) {
    int f = 0; cerr << '{';
    for (auto &i : x) cerr << (f++ ? "," : ""), __print(i);
    cerr << "}";
}
void _print() {cerr << "]\n";}
template <typename T, typename... V> void _print(T t, V... v) {
    __print(t); if (sizeof...(v)) cerr << ", "; _print(v...);
}
#ifndef ONLINE_JUDGE
#define debug(x...) cerr << "[" << #x << "] = ["; _print(x)
#else
#define debug(x...)
#endif
// freopen("input.txt", "r", stdin);
// freopen("output.txt", "w", stdout);
\end{lstlisting}
\subsection{Stress test}
\begin{lstlisting}
// compile the two solutions first:
//   g++ brute.cpp -o brute && g++ fast.cpp -o fast
#include <bits/stdc++.h>
using namespace std;
mt19937 rng(chrono::steady_clock::now().time_since_epoch().count());
long long rnd(long long l, long long r) {
    return uniform_int_distribution<long long>(l, r)(rng);
}
void gen() {
    ofstream in("in.txt");
    int n = rnd(1, 10); in << n << '\n';
    for(int i=0; i<n; i++) in << rnd(-100, 100) << " ";
}
int main() {
    for(int tc=1; ; tc++) {
        gen();
        system("./brute < in.txt > b.txt");
        system("./fast < in.txt > f.txt");
        if(system("diff -w b.txt f.txt")) {
            cout << "Mismatch on Test " << tc
                 << "! Check in.txt, b.txt, and f.txt\n";
            break;
        }
        cout << "Test " << tc << " OK\n";
    }
}
\end{lstlisting}
\subsection{VS Code setup}
\begin{lstlisting}
// keybindings.json - F5 compiles and runs in.txt -> out.txt
{ "key": "f5",
  "command": "workbench.action.terminal.sendSequence",
  "args": {
    "text": "clear && g++ ${fileBasenameNoExtension}.cpp -o ${fileBasenameNoExtension} && ./${fileBasenameNoExtension} < in.txt > out.txt\n"
  }
}
// settings.json also needs: "files.autoSaveDelay": 100
\end{lstlisting}
% Formula section - typeset by hand (the Word original kept these as images/OMML)
\section{Formula}

\subsection{Area}
\noindent\begin{tabular}{@{}l@{\hspace{1em}}l@{}}
Rectangle          & $\ell \times w$\\
Square             & $s \times s$\\
Triangle           & $\tfrac{1}{2}\,b\,h$\\
Parallelogram      & $b\,h$\\
Pyramid (no base)  & $\tfrac{1}{2}\,P_{\text{base}}\,\ell_{\text{slant}}$\\
\end{tabular}

\smallskip
\noindent Polygon, shoelace:
\[A=\tfrac{1}{2}\Bigl|\textstyle\sum_{i=1}^{n-1}\bigl(x_iy_{i+1}-x_{i+1}y_i\bigr)\Bigr|\]
\noindent Polygon, Pick (integer coordinates): $A=a+\frac{b}{2}-1$, where
$a$ = interior lattice points, $b$ = lattice points on the sides.

\subsection{Volume}
\noindent\begin{tabular}{@{}l@{\hspace{1em}}l@{}}
Cube               & $s^3$\\
Rectangular prism  & $\ell\,w\,h$\\
Cylinder           & $\pi r^2 h$\\
Sphere             & $\tfrac{4}{3}\pi r^3$\\
Pyramid            & $\tfrac{1}{3}\,A_{\text{base}}\,h$\\
\end{tabular}

\subsection{Surface Area}
\noindent\begin{tabular}{@{}l@{\hspace{1em}}l@{}}
Cube          & $6s^2$\\
Rect. prism   & $2(\ell w+\ell h+wh)$\\
Cylinder      & $2\pi r(r+h)$\\
Sphere        & $4\pi r^2$\\
Pyramid       & $A_{\text{base}}+\tfrac{1}{2}P_{\text{base}}\,\ell_{\text{slant}}$\\
\end{tabular}

\subsection{Triangles}
\noindent Side lengths $a$, $b$, $c$:
\[s=\frac{a+b+c}{2}\qquad A=\sqrt{s(s-a)(s-b)(s-c)}\]
\[R=\frac{abc}{4A}\qquad r=\frac{A}{s}\]
\[m_a=\tfrac{1}{2}\sqrt{2b^2+2c^2-a^2}\qquad
  t_a=\sqrt{bc\Bigl[1-\Bigl(\tfrac{a}{b+c}\Bigr)^{2}\Bigr]}\]
\noindent($m_a$ = median to $a$, $t_a$ = angle bisector from $A$.)

\subsection{Trigonometry}
\[\frac{\sin\alpha}{a}=\frac{\sin\beta}{b}=\frac{\sin\gamma}{c}=\frac{1}{2R}
  \qquad a^2=b^2+c^2-2bc\cos\alpha\]
\[\frac{a+b}{a-b}=\frac{\tan\frac{\alpha+\beta}{2}}{\tan\frac{\alpha-\beta}{2}}\]
\[\sin(A\pm B)=\sin A\cos B\pm\cos A\sin B\]
\[\cos(A\pm B)=\cos A\cos B\mp\sin A\sin B\]
\[\tan(A\pm B)=\frac{\tan A\pm\tan B}{1\mp\tan A\tan B}\]
\[\sin2\theta=2\sin\theta\cos\theta\qquad
  \cos2\theta=\cos^2\theta-\sin^2\theta\]
\[\tan2\theta=\frac{2\tan\theta}{1-\tan^2\theta}\]
\[\sin\tfrac{\theta}{2}=\pm\sqrt{\tfrac{1-\cos\theta}{2}}\qquad
  \cos\tfrac{\theta}{2}=\pm\sqrt{\tfrac{1+\cos\theta}{2}}\]
\[\tan\tfrac{\theta}{2}=\frac{1-\cos\theta}{\sin\theta}\]
\[\sin r+\sin w=2\sin\tfrac{r+w}{2}\cos\tfrac{r-w}{2}\]
\[\cos r+\cos w=2\cos\tfrac{r+w}{2}\cos\tfrac{r-w}{2}\]
\[(V+W)\tan\tfrac{r-w}{2}=(V-W)\tan\tfrac{r+w}{2}\]
\noindent where $V$, $W$ are the lengths of the opposite sides.
\[a\cos x+b\sin x=r\cos(x-\varphi)\qquad
  a\sin x-b\cos x=r\sin(x-\varphi)\]
\noindent with $r=\sqrt{a^2+b^2}$, $\varphi=\operatorname{atan2}(b,a)$.

\subsection{Sums}
\[c^k+c^{k+1}+\cdots+c^n=\frac{c^{n+1}-c^k}{c-1}\quad (c\neq 1)\]
\[\sum_{i=1}^{n} i=\frac{n(n+1)}{2}\qquad
  \sum_{i=1}^{n} i^2=\frac{n(n+1)(2n+1)}{6}\]
\[\sum_{i=1}^{n} i^3=\frac{n^2(n+1)^2}{4}\qquad
  \sum_{i=1}^{n} i^4=\frac{n(n+1)(2n+1)(3n^2+3n-1)}{30}\]
\noindent Sum of the first $n$ odd numbers $=n^2$.

\subsection{Logarithms}
\[\log_b 1=0\qquad \log_b b=1\qquad b^{\log_b a}=a\]
\[x^{\log_b y}=y^{\log_b x}\qquad \log_a b=\frac{1}{\log_b a}\]
\[\log_a x=\frac{\log_b x}{\log_b a}\qquad \log_a c=\log_a b\cdot\log_b c\]
\[\log_b(AB)=\log_b A+\log_b B\qquad
  \log_b\!\Bigl(\frac{A}{B}\Bigr)=\log_b A-\log_b B\]
\[\log_b A^x=x\log_b A\]

\subsection{Series}
\noindent\textbf{Catalan}\quad 1, 1, 2, 5, 14, 42, 132, 429, \dots
\[C_n=\frac{1}{n+1}\binom{2n}{n}=\frac{(2n)!}{(n+1)!\,n!}\quad (n\ge 0)\]
\[C_n=C_0C_{n-1}+C_1C_{n-2}+\cdots+C_{n-1}C_0\]
\noindent\textbf{Arithmetic}
\[a_n=a+(n-1)d\qquad s_n=\tfrac{n}{2}\bigl(2a+(n-1)d\bigr)\]
\noindent\textbf{Geometric}
\[a_n=a\,r^{\,n-1}\qquad s_n=\frac{a(1-r^n)}{1-r}\]
\noindent\textbf{Derangement}\quad 0, 1, 2, 9, 44, 265, 1854,
14833, 133496, 1334961, 14684570
\[D_n=n!\sum_{k=0}^{n}\frac{(-1)^k}{k!}\qquad
  D_n=\Bigl\lfloor\frac{n!}{e}+\frac{1}{2}\Bigr\rfloor\]
\noindent\textbf{$n$-th Fibonacci (golden ratio)}
\[f_n=\frac{1}{\sqrt5}\left[\Bigl(\frac{1+\sqrt5}{2}\Bigr)^{n}-
  \Bigl(\frac{1-\sqrt5}{2}\Bigr)^{n}\right]\]

\subsection{Facts}
\[\Bigl\lceil\frac{a}{b}\Bigr\rceil=\Bigl\lfloor\frac{a-1}{b}\Bigr\rfloor+1
  \qquad
  \sum_{i=l}^{r} i=\frac{l+r}{2}\,(r-l+1)\]
\[\Bigl\lfloor\frac{\lfloor n/a\rfloor}{b}\Bigr\rfloor
  =\Bigl\lfloor\frac{n}{ab}\Bigr\rfloor\]

\section{Number Theory}
\subsection{BINARY EXPONENTIATION:(a\^{}b)}
\begin{lstlisting}
long long binpow(long long a, long long b, long long m) {
    a %= m;
    long long res = 1;
    while (b > 0) {
        if (b & 1)
            res = res * a % m;
        a = a * a % m;
        b >>= 1;
    }
    return res;
}
\end{lstlisting}
\subsection{BINARY EXPONENTIATION:(a\^{}b\^{}c)}
\begin{lstlisting}
int binaryExp(int base, int power, int modulo){
    int ans = 1;
    while (power){
        if (power % 2 == 1)
            ans = (ans * base) % modulo;
        base = (base * base) % modulo;
        power /= 2;
    }
    return ans;
} //function call:
binaryExp(a, binaryExp(b, c, mod-1), mod)
\end{lstlisting}
\subsection{Prime factorization-O(sqrt(n))}
\[\sigma(N)=\frac{p_1^{a+1}-1}{p_1-1}\cdot\frac{p_2^{b+1}-1}{p_2-1}\cdot\frac{p_3^{c+1}-1}{p_3-1}\]
\begin{lstlisting}
// smallest prime factor of a number.
int factor(int n){
    int a;
    if (n%2==0)
        return 2;
    for (a=3; a<=sqrt(n); a+=2){
        if (n%a==0)
            return a;
    }
    return n;
}
// complete factorization
int r;
while (n>1){
    r = factor(n);
    printf("%d", r);
    n /= r;
}
\end{lstlisting}
\subsection{smallest prime factor(SPF) 1 to N}
\begin{lstlisting}
const int N = 1e7 + 5;
int spf[N];
void smallestPrimeFactorUsingSeive() {
    for (int i = 2; i < N; i++) {
        if (spf[i] == 0) {
            for (int j = i; j < N; j += i) {
                if (spf[j] == 0)
                    spf[j] = i;
            }
        }
    }
}
\end{lstlisting}
\subsection{Sieve}
\begin{lstlisting}
const ll N = 1e7 + 5;
ll isprime[N];
vector<int> primes;
void sieveOfEratosthenes() {
    for (ll i = 2; i < N; i++)
        isprime[i] = 1;
    for (ll i = 4; i < N; i += 2)
        isprime[i] = 0;
    for (ll i = 3; i * i < N; i += 2) {
        if (isprime[i]) {
            for (ll j = i * i; j<N; j+=i*2)
                isprime[j] = 0;
        }
    }
    for (ll i = 2; i < N; i++)
        if (isprime[i])
            primes.push_back(i);
}
\end{lstlisting}
\subsection{Segmented Sieve}
\begin{lstlisting}
// find all the prime num in a range [L,R] of small size (R-L+1==1e7) where R can be very large e.g 1e12.
// before that you need to call sieve and store the primes up to sqrt(R)
void segmentedSieve(ll L, ll R) {
    bool isPrime[R - L + 1];
    for (ll i = 0; i <= R - L; i++)
        isPrime[i] = true;
    if (L == 1)
        isPrime[0] = false;
    for (ll i=0;primes[i]*primes[i]<=R;i++){
        ll curPrime = primes[i];
        ll base = curPrime * curPrime;
        if (base < L) {
            base = ((L + curPrime - 1) / curPrime) * curPrime;
        }
        for (ll j=base;j<=R; j += curPrime)
            isPrime[j - L] = false;
    }
    for (ll i = 0; i <= R - L; i++) {
        if (isPrime[i] == true)
            cout << L + i << " ";
    }
    cout << endl;
}
\end{lstlisting}
\subsection{Bitwise Sieve}
\begin{lstlisting}
const ll N = 10000006;
bitset<N> sieve;
void bitwiseSieve() {
    sieve.flip();
    ll finalBit = sqrt(sieve.size()) + 1;
    for (ll i = 2; i < finalBit; i++) {
        if (sieve.test(i))
            for (ll j = 2 * i; j<N; j += i)
                sieve.reset(j);
    }
}
// to check any number is prime or not,
// check sieve.test(number) is true or not
\end{lstlisting}
\subsection{Sieve up to 1e9 in 500ms}
\begin{lstlisting}
// takes 0.5s for n = 1e9
// copy from YouKnowWho
vector<ll> sieve(ll N, ll Q=17,ll L=1<<15){
    ll rs[] = {1,7,11,13, 17, 19, 23, 29};
    struct P {
        P(ll p) : p(p) {}
        ll p;
        ll pos[8];
    };
    auto approx_prime_count = [](ll N) {
        return N > 60184 ? N / (log(N)-1.1)
        : max(1., N / (log(N) - 1.11)) + 1;
    };
    ll v = sqrt(N), vv = sqrt(v);
    vector<bool> isp(v + 1, true);
    for (ll i = 2; i <= vv; ++i) {
        if (isp[i]) {
            for (ll j = i*i;j<=v;j+=i)
                isp[j] = false;
        }
    }
    ll rsize = approx_prime_count(N + 30);
    vector<ll> primes = {2, 3, 5};
    ll ps = 3;
    primes.resize(rsize);
    vector<P> sprimes;
    size_t pbeg = 0;
    ll prod = 1;
    for (ll p = 7; p <= v; ++p) {
        if (!isp[p])
            continue;
        if (p <= Q) {
            prod *= p;
            ++pbeg;
            primes[ps++] = p;
        }
        auto pp = P(p);
        for (ll t = 0; t < 8; ++t) {
            ll j = (p <= Q) ? p : p * p;
            while (j % 30 != rs[t])
                j += p << 1;
            pp.pos[t] = j / 30;
        }
        sprimes.push_back(pp);
    }
    vector<unsigned char> pre(prod, 0xFF);
    for (size_t pi = 0; pi < pbeg; ++pi) {
        auto pp = sprimes[pi];
        ll p = pp.p;
        for (ll t = 0; t < 8; ++t) {
            unsigned char m = ~(1 << t);
            for (ll i=pp.pos[t];i<prod;i+=p)
                pre[i] &= m;
        }
    }
    ll bs = (L + prod - 1) / prod * prod;
    vector<unsigned char> block(bs);
    unsigned char* pb = block.data();
    ll M = (N + 29) / 30;
    for (ll s = 0; s < M; s += bs,pb-=bs) {
        ll f = min(M, s + bs);
        for (ll i = s; i < f; i += prod) {
        copy(pre.begin(), pre.end(), pb+i);
        }
        if (s == 0)  pb[0] &= 0xFE;
        for (size_t pi = pbeg; pi < sprimes.size(); ++pi) {
            auto& pp = sprimes[pi];
            ll p = pp.p;
            for (ll t = 0; t < 8; ++t) {
                ll i = pp.pos[t];
                unsigned char m = ~(1 << t);
                for (; i < f; i += p)
                    pb[i] &= m;
                pp.pos[t] = i;
            }
        }
        for (ll i = s; i < f; ++i) {
            for (ll m=pb[i];m>0;m&=m-1){
     primes[ps++]=i*30+rs[__builtin_ctz(m)];
            }
        }
    }
    assert(ps <= rsize);
    while (ps > 0 && primes[ps - 1] > N)
        --ps;
    primes.resize(ps);
    return primes;
}
\end{lstlisting}
\subsection{Legendre formula(Maximum Power of n!)}
\begin{lstlisting}
// maximum power of prime p that divides n!
int legendre(int n, int p) {
    int ans = 0;
    while (n) {
        n /= p;
        ans += n;
    }
    return ans;
}
\end{lstlisting}
\subsection{EXT\_GCD}
\begin{lstlisting}
// return {x,y} such that ax+by=gcd(a,b)
pair<int,int>ext_gcd(int a, int b){
    if (b == 0)
        return {1, 0};
    else{
        pair<int,int> tmp=ext_gcd(b, a % b);
        return {tmp.second, tmp.first - (a / b) * tmp.second};
    }
}
\end{lstlisting}
\subsection{PHI of N}
\[n=p_1^{a_1}p_2^{a_2}\cdots p_k^{a_k}\ \Rightarrow\ \varphi(n)=n\Bigl(1-\tfrac{1}{p_1}\Bigr)\Bigl(1-\tfrac{1}{p_2}\Bigr)\cdots\Bigl(1-\tfrac{1}{p_k}\Bigr)\]
\begin{lstlisting}
// the positive integers less than or equal to n that are relatively prime to n.
int phi(int n) {
    int result = n;
    for (int i = 2; i * i <= n; i++) {
        if (n % i == 0) {
            while (n % i == 0)
                n /= i;
            result -= result / i;
        }
    }
    if (n > 1)
        result -= result / n;
    return result;
}
\end{lstlisting}
\subsection{PHI of 1 to N}
\[\varphi(1)+\varphi(2)+\varphi(5)+\varphi(10)=1+1+4+4=10\]
\begin{lstlisting}
const int N = 1e5 + 5;
vector<int> phi(N);
void phi_1_to_n() {
    for (int i = 0; i < N; i++)
        phi[i] = i;
    for (int i = 2; i < N; i++) {
        if (phi[i] == i) {
            for (int j = i; j < N; j += i)
                phi[j] -= phi[j] / i;
        }
    }
}
// Fact: the phi of every divisor of N sums to N.
// N = 10: divisors 1, 2, 5, 10 -> 1+1+4+4 = 10
\end{lstlisting}
\subsection{nCr(more space, less time)}
\begin{lstlisting}
int mod = 1e9 + 7;
const int MAX = 1e7 + 5;
vector<int> fact(MAX), ifact(MAX), inv(MAX);
void factorial() {
    inv[1] = fact[0] = ifact[0] = 1;
    for (int i = 2; i < MAX; i++)
        inv[i]=inv[mod%i]*(mod-mod/i)%mod;
    for (int i = 1; i < MAX; i++)
        fact[i] = (fact[i - 1] * i) % mod;
    for (int i = 1; i < MAX; i++)
        ifact[i]=ifact[i-1]*inv[i] % mod;
}

int nCr(int n, int r) {
    if (r < 0 || r > n)
        return 0;
    return (int)fact[n] * ifact[r] % mod * ifact[n - r] % mod;
}
// first call factorial() function
// then for nCr just call nCr(n,r)
\end{lstlisting}
\subsection{nCr (less space, more time)}
\begin{lstlisting}
const int MOD = 1e9 + 7;
const int MAX = 1e7+10;
vector<int> fact(MAX), inv(MAX);
void factorial(){
    fact[0] = 1;
    for (int i = 1; i < MAX; i++)
        fact[i] = (i * fact[i - 1]) % MOD;
}
//For binaryExp we call 1.6 function
void inverse(){
    for (int i = 0; i < MAX; ++i)
        inv[i]=binaryExp(fact[i], MOD - 2, MOD);
}
int nCr(int a, int b){
    if (a < b or a < 0 or b < 0)
        return 0;
    int de = (inv[b] * inv[a - b]) % MOD;
    return (fact[a] * de) % MOD;
}
// nCr ends here
int ModInv(int a, int M){
    return binaryExp(a, M - 2, M);
}
\end{lstlisting}
\subsection{Factorial mod}
\begin{lstlisting}
//n! mod p : Here P is mod value
//For binaryExp we call 1.6 function
int factmod (int n, int p) {
    int res = 1;
    while (n > 1){
        res=(res*binaryExp(p-1,n/p,p))%p;
        for (int i=2; i<=n%p; ++i)
            res=(res*i) %p;
            n /= p;
    }
    return int (res % p);
}
\end{lstlisting}
\subsection{Modular Operation}
\begin{lstlisting}
// Addition
int mod_add(int a, int b, int MOD = mod){
    a = a % MOD, b = b % MOD;
    return (((a + b) % MOD) + MOD) % MOD;
}
// Subtraction
int mod_sub(int a, int b, int MOD = mod){
    a = a % MOD, b = b % MOD;
    return (((a - b) % MOD) + MOD) % MOD;
}
// Multiplication
int mod_mul(int a, int b, int MOD = mod){
    a = a % MOD, b = b % MOD;
    return (((a * b) % MOD) + MOD) % MOD;
}
// Division
//call binary Exponential Function here.
int mminvprime(int a, int b) { return binaryExp(a, b - 2, b); }
//call modular multiplication here.
int mod_div(int a, int b, int MOD = mod) {
    a = a % MOD, b = b % MOD;
    return (mod_mul(a, mminvprime(b, MOD), MOD) + MOD) % MOD;
}
//only for prime MOD
\end{lstlisting}
\subsection{Convex Hull}
\begin{lstlisting}
#define pld pair<ld, ld>
ll orientation(pld a, pld b, pld c) {
//Cal the determinant to determine orientat
    double v = a.first*(b.second-c.second)+
               b.first*(c.second-a.second)+
               c.first*(a.second-b.second);
    return (v < 0) ? -1 : (v > 0) ? 1 : 0;  // -1: clockwise, 1: counter-clockwise, 0: collinear
}
bool cw(pld a,pld b,pld c, bool collinear) {
// Check if points (a, b, c) are clockwise or collinear (if allowed)
    ll o = orientation(a, b, c);
    return o < 0 || (collinear && o == 0);
}

bool ccw(pld a, pld b,pld c,bool collinear){
// Check if points (a,b,c) are counter-clockwise
// or collinear (if allowed)
    ll o = orientation(a, b, c);
    return o > 0 || (collinear && o == 0);
}

void convex_hull(vector<pld> &a, bool collinear = false) {
    if (a.size() == 1)
        return;  // single point, no hull needed
// Step 1: Sort points by x-coordinate, then by y-coordinate
    sort(a.begin(), a.end(), [](pld a, pld b) {
        return make_pair(a.first, a.second) <
               make_pair(b.first, b.second);
    });
    pld p1 = a[0], p2 = a.back();
// Leftmost and rightmost points
    vector<pld> up = {p1}, down = {p1};
// Upper and lower hulls
    for (ll i = 1; i < (int)a.size(); i++) {
// Add to upper hull if last point or forms clockwise turn
        if (i == a.size() - 1 ||
cw(p1, a[i], p2, collinear)) {
            while (up.size() >= 2 && !cw(up[up.size() - 2], up[up.size() - 1], a[i], collinear)) { up.pop_back(); }
            up.push_back(a[i]);
        }
// Add to lower hull if last point or forms counter-clockwise turn
        if (i == a.size() - 1 ||
ccw(p1, a[i], p2, collinear)) {
            while (down.size() >= 2 && !ccw(down[down.size() - 2], down[down.size() - 1], a[i], collinear))
               { down.pop_back(); }
            down.push_back(a[i]);
        }
    }
// Handle special case:all points are collinear
    if (collinear && up.size()==a.size()) {
        reverse(a.begin(), a.end());
        return;
    }
// Merge upper and lower hulls into the result
    a.clear();
    for (ll i = 0; i < (int)up.size(); i++)
        a.push_back(up[i]);
    for (ll i = down.size() - 2; i > 0; i--)
        a.push_back(down[i]);
}
\end{lstlisting}
\subsection{Line Intersection}
\begin{lstlisting}
int ori(const pair<int, int> &o, const pair<int, int> &a, const pair<int, int> &b)
{
    int ret = (a.first - o.first) * (b.second - o.second) -
              (a.second - o.second) * (b.first - o.first);
    return (ret > 0) - (ret < 0);
}

bool isIntersect(const pair<int, int> &p1, const pair<int, int> &p2,
                 const pair<int, int> &q1, const pair<int, int> &q2)
{
    if (ori(p1, p2, q1) == 0 && ori(p1, p2, q2) == 0 && ori(q1, q2, p1) == 0 && ori(q1, q2, p2) == 0)
    {
        if (ori(p1, p2, q1))
            return false;
        return (p1.first - q1.first) * (p2.second - q1.second) <= 0 ||
               (p1.second - q1.second) * (p2.first - q1.first) <= 0 ||
               (q1.first - p1.first) * (q2.second - p1.second) <= 0 ||
               (q1.second - p1.second) * (q2.first - p1.first) <= 0;
    }
    return (ori(p1, p2, q1) * ori(p1, p2, q2) <= 0) &&
           (ori(q1, q2, p1) * ori(q1, q2, p2) <= 0);
}
\end{lstlisting}
\section{Algorithms}
\subsection{KMP Algorithm-O(n+m)}
\begin{lstlisting}
vector<int> createLPS(string pattern) {
    int n = pattern.length(), idx = 0;
    vector<int> lps(n);
    for (int i = 1; i < n;) {
        if (pattern[idx] == pattern[i]) {
            lps[i] = idx + 1;
            idx++, i++;
        }
        else {
            if (idx != 0)
                idx = lps[idx - 1];
            else
                lps[i] = idx, i++;
        }
    }
    return lps;
}
int kmp(string text, string pattern) {
    int cnt_of_match = 0, i = 0, j = 0;
    vector<int> lps = createLPS(pattern);
    while (i < text.length()) {
        if (text[i] == pattern[j])
            i++, j++;// i->text,j->pattern
        else {
            if (j != 0)
                j = lps[j - 1];
            else
                i++;
        }
        if (j == pattern.length()) {
            cnt_of_match++;
            // the index where match found -> (i - pattern.length());
            j = lps[j - 1];
        }
    }
    return cnt_of_match;
}
\end{lstlisting}
\subsection{2D prefix sum}
\begin{lstlisting}
class NumMatrix {
    int row, col;
    vector<vector<int>> sums;
public:
    NumMatrix(vector<vector<int>> &matrix) {
        row = matrix.size();
        col = row>0 ? matrix[0].size() : 0;
        sums = vector<vector<int>>(row+1, vector<int>(col+1, 0));
        for(int i=1; i<=row; i++) {
            for(int j=1; j<=col; j++) {
                sums[i][j] = matrix[i-1][j-1] + sums[i-1][j] +
                             sums[i][j-1] - sums[i-1][j-1];
            }
        }
    }
    int sumRegion(int row1, int col1, int row2, int col2) {
        return sums[row2+1][col2+1] - sums[row2+1][col1] -
               sums[row1][col2+1] + sums[row1][col1];
    }
};
\end{lstlisting}
\subsection{Biginteger Operation}
\begin{lstlisting}
struct BigInteger {
    string str;
    // Constructor to initialize
    // BigInteger with a string
    BigInteger(string s) { str = s; }
    // Overload + operator to add
    // two BigInteger objects
    BigInteger operator+(const BigInteger& b) {
        string a = str, c = b.str;
        int alen=a.length(),clen=c.length();
        int n = max(alen, clen);
        if (alen > clen)
            c.insert(0, alen - clen, '0');
        else if (alen < clen)
            a.insert(0, clen - alen, '0');
        string res(n + 1, '0');
        int carry = 0;
        for (int i = n - 1; i >= 0; i--) {
            int digit = (a[i]-'0') + (c[i]-'0') + carry;
            carry = digit / 10;
            res[i + 1] = digit % 10 + '0';
        }
        if (carry == 1) {
            res[0] = '1';
            return BigInteger(res);
        }
        else
           return BigInteger(res.substr(1));
    }

    // Overload - operator to subtract
    // first check which number is greater and then subtract
    BigInteger operator-(const BigInteger& b) {
        string a = str;
        string c = b.str;
        int alen=a.length(),clen=c.length();
        int n = max(alen, clen);
        if (alen > clen)
            c.insert(0, alen - clen, '0');
        else if (alen < clen)
            a.insert(0, clen - alen, '0');
        if (a < c) {
            swap(a, c);
            swap(alen, clen);
        }
        string res(n, '0');
        int carry = 0;
        for (int i = n - 1; i >= 0; i--) {
            int digit = (a[i]-'0') - (c[i]-'0') - carry;
            if (digit < 0) {
                digit += 10;
                carry = 1;
            }
            else {
                carry = 0;
            }
            res[i] = digit + '0';
        }
        // remove leading zeros
        int i = 0;
        while (i < n && res[i] == '0')
            i++;
        if (i == n)
            return BigInteger("0");
        return BigInteger(res.substr(i));
    }

    // Overload * operator to multiply
    // two BigInteger objects
    BigInteger operator*(const BigInteger& b) {
        string a = str, c = b.str;
        int alen=a.length(),clen=c.length();
        int n = alen + clen;
        string res(n, '0');
        for (int i = alen - 1; i >= 0;i--) {
            int carry = 0;
            for(int j=clen-1; j>=0; j--) {
                int digit = (a[i]-'0') * (c[j]-'0') +
                            (res[i+j+1]-'0') + carry;
                carry = digit / 10;
                res[i+j+1]=digit % 10 + '0';
            }
            res[i] += carry;
        }
        int i = 0;
        while (i < n && res[i] == '0')
            i++;
        if (i == n)
            return BigInteger("0");
        return BigInteger(res.substr(i));
    }
    // Overload << operator to output
    // BigInteger object
    friend ostream& operator<<(ostream& out, const BigInteger& b) {
        out << b.str;
        return out;
    }
};
\end{lstlisting}
\subsection{Grundy Number}
\begin{lstlisting}
int mex(unordered_set<int> s) {
    int m = 0;
    while (s.count(m)) m++;
    return m;
}

int grundy(int n, vector<int>& moves, vector<int>& dp) {
    if (n == 0) return 0;
    if (dp[n] != -1) return dp[n];

    unordered_set<int> s;
    for (int move : moves) {
        if (n >= move)
            s.insert(grundy(n - move, moves, dp));
    }
    return dp[n] = mex(s);
}

int main() {
    vector<int> moves = {1, 2, 3}; // allowed moves
    int n = 10; // number of stones
    vector<int> dp(n + 1, -1);

    for (int i = 0; i <= n; i++)
        cout << "G(" << i << ") = " << grundy(i, moves, dp) << endl;

    return 0;
}
\end{lstlisting}
\subsection{Manacher}
\begin{lstlisting}
struct Manacher {
  vector<int> p[2];
  // p[1][i] = (max odd length palindrome centered at i) / 2 [floor division]
  // p[0][i] = same for even, it considers the right center
  // e.g. for s = "abbabba", p[1][3] = 3, p[0][2] = 2
  Manacher(string s) {
    int n = s.size();
    p[0].resize(n + 1);
    p[1].resize(n);
    for (int z = 0; z < 2; z++) {
      for (int i = 0, l = 0, r = 0; i < n; i++) {
        int t = r - i + !z;
        if (i < r) p[z][i] = min(t, p[z][l + t]);
        int L = i - p[z][i], R = i + p[z][i] - !z;
        while (L >= 1 && R + 1 < n && s[L - 1] == s[R + 1])
          p[z][i]++, L--, R++;
        if (R > r) l = L, r = R;
      }
    }
  }
  bool is_palindrome(int l, int r) {
    int mid = (l + r + 1) / 2, len = r - l + 1;
    return 2 * p[len % 2][mid] + len % 2 >= len;
  }
};
\end{lstlisting}
\subsection{Misc snippets}
\begin{lstlisting}
// binary search on the answer
bool good(int w, int n, int c, vector<int>& a);
int l = 0;              // try l = -1 if not ac
int r = 2;
while (!good(r, n, c, a)) r *= 2;
while (r > l + 1) {
    int m = (l + r) / 2;
    if (good(m, n, c, a)) r = m;
    else l = m;
}
cout << r << '\n';

do {
        for (auto i: v) cout << i << " ";
        cout << endl;
    } while (next_permutation(v.begin(), v.end()));
#define multiply(a,b) (((a)%mod)*((b)%mod))%mod
#define fr(i,a,b) for (int i = a; i <= b; i++)
#define all(v) (v).begin(), (v).end()
#define endl '\n'

vector<vector<int>> arr(rows, vector<int>(cols, 0));
size_t pos = s.find("w"); //string
if (pos != string::npos) found;
binpow(a-1,mod-2,mod) //division
__builtin_popcountll(n)// 0 not allowed
// also __builtin_ctzll, __builtin_clzll
cout << fixed << setprecision(15) << answer << "\n"; //use long double safe for 21 places

// misere nim: special case when all piles are 1

// time ./a.out < input.txt

v.erase(unique(v.begin(), v.end()), v.end());
\end{lstlisting}
\section{Data Structure}
\subsection{SEGMENT TREE}
\begin{lstlisting}
class SEGMENT_TREE {
   public:
    vector<int> v;
    vector<int> seg;
    SEGMENT_TREE(int n) {
        v.resize(n + 5);
        seg.resize(4 * n + 5);
    }
    //! initially: ti = 1, low = 1, high = n
//(number of elements in the array);
    void build(int ti, int low, int high) {
        if (low == high) {
            seg[ti] = v[low];
            return;
        }
        int mid = (low + high) / 2;
        build(2 * ti, low, mid);
        build(2 * ti + 1, mid + 1, high);
        seg[ti] = (seg[2*ti]+seg[2*ti+1]);
    }
    //! initially: ti = 1, low = 1, high = n
//(number of elements in the array),
//(ql & qr)=user input in 1 based index;
    int find(int ti, int tl, int tr, int ql, int qr) {
        if (tl > qr || tr < ql) {
            return 0;
        }
        if (tl >= ql and tr <= qr)
            return seg[ti];
        int mid = (tl + tr) / 2;
        int l = find(2*ti, tl, mid, ql, qr);
        int r = find(2*ti+1,mid+1,tr,ql,qr);
        return (l + r);
    }
    //! initially: ti = 1, tl = 1, tr = n
//(number of elements in the array),
//id = user input in 1 based indexing,
//val = updated value;
    void update(int ti, int tl, int tr, int id, int val) {
        if (id > tr or id < tl)
            return;
        if (id == tr and id == tl) {
            seg[ti] = val;
            return;
        }
        int mid = (tl + tr) / 2;
        update(2 * ti, tl, mid, id, val);
        update(2*ti+1,mid + 1, tr, id, val);
        seg[ti] = (seg[2*ti]+seg[2*ti + 1]);
    }
};
// use 1 based indexing for input and //queries and update;
\end{lstlisting}
\subsection{SEGMENT TREE LAZY}
\begin{lstlisting}
const int N = 1e5 + 100;
int tree[N << 2], lz[N << 2];
void propagate(int u, int st, int en) {
    if (!lz[u])
        return;
    tree[u] += lz[u] * (en - st + 1);
    if (st != en) {
        lz[2 * u] += lz[u];
        lz[2 * u + 1] += lz[u];
    }
    lz[u] = 0;
}
void update(int u, int st, int en, int l, int r, int x) {
    propagate(u, st, en);
    if (r < st or en < l)
        return;
    else if (st >= l and en <= r) {
        lz[u] += x;
        propagate(u, st, en);
    }
    else {
        int mid = (st + en) >> 1;
        update(2 * u, st, mid, l, r, x);
        update(2*u + 1, mid+1, en, l, r, x);
        tree[u] = tree[2*u]+tree[2*u+1];
    }
}
int query(int u,int st,int en,int l,int r){
    propagate(u, st, en);
    if (r < st or en < l)
        return 0;
    else if (st >= l and en <= r)
        return tree[u];
    else {
        int mid = (st + en) >> 1;
        int left=query(2*u, st, mid, l, r);
        int right=query(2*u+1,mid+1,en,l,r);
        return left + right;
    }
}
\end{lstlisting}
\subsection{DSU}
\begin{lstlisting}
class DisjointSet{
    vector<int> par, sz, minElmt, maxElmt, cntElmt;

   public:
    DisjointSet(int n){
        par.resize(n + 1);
        sz.resize(n + 1, 1);
        minElmt.resize(n + 1);
        maxElmt.resize(n + 1);
        cntElmt.resize(n + 1, 1);
        for (int i = 1; i <= n; i++)
            par[i]=minElmt[i]=maxElmt[i]=i;
    }
    int findUPar(int u) {
        if (u == par[u])
            return u;
        return par[u] = findUPar(par[u]);
    }
    void unionBySize(int u, int v){
        int pU = findUPar(u);
        int pV = findUPar(v);
        if (pU == pV)
            return;
        if (sz[pU] < sz[pV])
            swap(pU, pV);
        par[pV] = pU;
        sz[pU] += sz[pV];
        cntElmt[pU] += cntElmt[pV];
        minElmt[pU] = min(minElmt[pU], minElmt[pV]);
        maxElmt[pU] = max(maxElmt[pU], maxElmt[pV]);
    }
    int getMinElementIntheSet(int u){
        return minElmt[findUPar(u)];
    }
    int getMaxElementIntheSet(int u){
        return maxElmt[findUPar(u)];
    }
    int getNumofElementIntheSet(int u){
        return cntElmt[findUPar(u)];
    }
};
\end{lstlisting}
\subsection{String Hashing}
\begin{lstlisting}
// include binary exponential here.
const ll N = 2e5 + 5;
const ll MOD1 = 127657753, MOD2 = 987654319;
const ll p1 = 137, p2 = 277;
ll ip1, ip2;
pair<ll, ll> pw[N], ipw[N];
void prec() {
    pw[0] = {1, 1};
    for (ll i = 1; i < N; i++) {
        pw[i].first = 1LL * pw[i-1].first * p1 % MOD1;
        pw[i].second = 1LL * pw[i-1].second * p2 % MOD2;
    }
    ip1 = binaryExp(p1, MOD1 - 2, MOD1);
    ip2 = binaryExp(p2, MOD2 - 2, MOD2);
    ipw[0] = {1, 1};
    for (ll i = 1; i < N; i++) {
        ipw[i].first = 1LL * ipw[i-1].first * ip1 % MOD1;
        ipw[i].second = 1LL * ipw[i-1].second * ip2 % MOD2;
    }
}
struct Hashing {
    ll n;
    string s;                 // 0 - indexed
    vector<pair<ll, ll>> hs;  // 1 - indexed
    Hashing() {}
    Hashing(string _s) {
        n = _s.size();
        s = _s;
        hs.emplace_back(0, 0);
        for (ll i = 0; i < n; i++) {
            pair<ll, ll> p;
            p.first = (hs[i].first + 1LL * pw[i].first * s[i] % MOD1) % MOD1;
            p.second = (hs[i].second + 1LL * pw[i].second * s[i] % MOD2) % MOD2;
            hs.push_back(p);
        }
    }
    pair<ll, ll> get_hash(ll l, ll r) {   // 1 - indexed
        assert(1 <= l && l <= r && r <= n);
        pair<ll, ll> ans;
        ans.first = (hs[r].first - hs[l - 1].first + MOD1) * 1LL * ipw[l - 1].first % MOD1;
        ans.second = (hs[r].second - hs[l - 1].second + MOD2) * 1LL * ipw[l - 1].second % MOD2;
        return ans;
    }
    pair<ll, ll> get_hash() { return get_hash(1, n); }
};
\end{lstlisting}
\subsection{Trie}
\begin{lstlisting}
const int N = 26;
char BASE = 'A';
class Node {
public:
    int EoW;
    Node *child[N];
    Node() {
        EoW = 0;
        for (int i = 0; i < N; i++)
            child[i] = NULL;
    }
};
Node *root = new Node();
void insert(Node *node, string s) {
    for (size_t i = 0; i < s.size(); i++) {
        int r = s[i] - BASE;
        if (node->child[r] == NULL)
            node->child[r] = new Node();
        node = node->child[r];
    }
    node->EoW += 1;
} // insert(root, s);
int search(Node *node, string s) {
    for (size_t i = 0; i < s.size(); i++) {
        int r = s[i] - BASE;
        if (node->child[r] == NULL)
            return 0;
        node = node->child[r];
    }
    return node->EoW;
} // search(root, s);
void print(Node *node, string s = "") {
    if (node->EoW)
        cout << s << "\n";
    for (int i = 0; i < N; i++) {
        if (node->child[i] != NULL) {
            char c = i + BASE;
            print(node->child[i], s + c);
        }
    }
} // print whole trie
void remove(Node *node, string s) {
    for (size_t i = 0; i < s.size(); i++) {
        int r = s[i] - BASE;
        if (node->child[r] == NULL)
            return;
        node = node->child[r];
    }
    node->EoW--;
} // remove(root, s)
void delete_trie(Node *node) {
    for (int i = 0; i < N; i++) {
        if (node->child[i] != NULL)
            delete_trie(node->child[i]);
    }
    delete node;
} // delete full trie
\end{lstlisting}
\subsection{Order Set}
\begin{lstlisting}
#include <ext/pb_ds/assoc_container.hpp>
#include <ext/pb_ds/tree_policy.hpp>
using namespace __gnu_pbds;

template <typename T> using o_set = tree<T, null_type, less<T>, rb_tree_tag, tree_order_statistics_node_update>;

// find_by_order(k) - returns an iterator to
// the k-th largest element (0 indexed);
// order_of_key(k)-the number of elements in // the set that are strictly smaller than k;
\end{lstlisting}
\subsection{GP Hash Table}
\begin{lstlisting}
#include <ext/pb_ds/assoc_container.hpp>
#include <ext/pb_ds/tree_policy.hpp>
using namespace __gnu_pbds;
template <typename p, typename q>
using ht = gp_hash_table<p, q>;
\end{lstlisting}
\section{Dynamic Programming}
\subsection{LCS O(n*m)}
\begin{lstlisting}
string s = "abbced", t = "bedc";
cin >> s >> t;
int n = s.size(), m = t.size();
int dp[n + 5][m + 5];
memset(dp, 0, sizeof(dp));
for (int i = 1; i <= n; i++) {
    for (int j = 1; j <= m; j++) {
        if (s[i - 1] == t[j - 1]) {
            dp[i][j] = 1 + dp[i - 1][j - 1];
        }
        else {
            dp[i][j] = max(dp[i-1][j], dp[i][j-1]);
        }
    }
}

// Problems that reduce to LCS:
// Longest Increasing Substring: only the
//   equal-char case matters, drop the rest.
// LPS: run LCS of s against reverse(s).
// Min insertions to make s a palindrome:
//   s.length() - LPS.  s = "aabca", LPS =
//   "aca": insert reverse("ab") after c ->
//   "aabcbaa".
// Min insert+delete to make s and t equal:
//   n + m - 2 * LCS.length()
\end{lstlisting}
\subsection{MCM O(n\^{}3)}
\begin{lstlisting}
const int N = 1005;
vector<int> v;
int dp[N][N], mark[N][N];
int MCM(int i, int j) {
    if (i == j)
        return dp[i][j] = 0;
    if (dp[i][j] != -1)
        return dp[i][j];
    int mn = INT_MAX;
    for (int k = i; k < j; k++) {
        int x = mn;
        mn = min(mn, MCM(i, k) + MCM(k + 1, j) + v[i - 1] * v[k] * v[j]);
        if (x != mn)
            mark[i][j] = k;
    }
    return dp[i][j] = mn;
}

void print_order(int i, int j) {
    if (i == j)
        cout << "X" << i;
    else {
        cout << "(";
        print_order(i, mark[i][j]);
        print_order(mark[i][j] + 1, j);
        cout << ")";
    }
}
// memset(dp, -1, sizeof dp);
// print_order(1, n);
\end{lstlisting}
\subsection{Length of LIS O(nlogn)}
\begin{lstlisting}
vector<int> v = {7, 3, 5, 3, 6, 2, 9, 8};
vector<int> seq;
/*
here we basically check is the current element from v is greater than the last
element of the sequence.
 if it is then push it to the seq array and if not then replace that index value.
let's take an example: v = 7 3 5 3 6 2 9 8 1st iteration seq = 7;
2nd iteration seq = 3;
3rd iteration seq = 3 5;
4th iteration seq = 3 3;
5th iteration seq = 3 3 6;
6th iteration seq = 2 3 6;
7th iteration seq = 2 3 6 9;
8th iteration seq = 2 3 6 8;
*/
for (auto i : v) {
    auto id = lower_bound(seq.begin(), seq.end(), i);
    if (id == seq.end())
        seq.push_back(i);
    else
        seq[id - seq.begin()] = i;
}
cout << seq.size() << endl;
\end{lstlisting}
\subsection{LCIS O(n * m)}
\begin{lstlisting}
int a[100]= {0}, b[100]= {0}, f[100]= {0};
int n=0, m=0;
int main(void){
    cin >> n;
    for (int i=1; i<=n; i++) cin >> a[i];
    cin >> m;
    for (int i=1; i<=m; i++) cin >> b[i];
    for (int i=1; i<=n; i++){
        int k=0;
        for (int j=1; j<=m; j++){
            if (a[i]>b[j] && f[j]>k)
k=f[j];
            else if (a[i]==b[j] && k+1>f[j])
f[j]=k+1;
        }
    }
    int ans=0;
    for (int i=1; i<=m; i++)
        if (f[i]>ans) ans=f[i];
    cout << ans << endl;
    return 0;
}
\end{lstlisting}
\subsection{Maximum submatrix}
\begin{lstlisting}
int a[150][150]= {0};
int c[200]= {0};
int maxarray(int n){
    int b=0, sum=-100000000;
    for (int i=1; i<=n; i++){
        if (b>0) b+=c[i];
        else b=c[i];
        if (b>sum) sum=b;
    }
    return sum;
}

int maxmatrix(int n){
    int sum=-100000000, max=0;
    for (int i=1; i<=n; i++){
        for (int j=1; j<=n; j++)
            c[j]=0;
        for (int j=i; j<=n; j++){
            for (int k=1; k<=n; k++)
                c[k]+=a[j][k];
            max=maxarray(n);
            if (max>sum) sum=max;
        }
    }
    return sum;
}
int main(void){
    int n=0;
    cin >> n;
    for (int i=1; i<=n; i++)
        for (int j=1; j<=n; j++)
            cin >> a[i][j];
    cout << maxmatrix(n);
    return 0;
}
\end{lstlisting}
\subsection{SOS DP}
\begin{lstlisting}
// # of elements in the list for which you // want to find the sum over all subsets
int n = 20;
// the list for which you want to find the // sum over all subsets
vector<int> a(1 << n);

//answer for sum over subsets of each subset
vector<int> sos(1 << n);

// O(4^n) brute force: check every j against every i
for (int i = 0; i < (1 << n); i++)
    for (int j = 0; j < (1 << n); j++)
        if ((i & j) == j)
            sos[i] += a[j];

// the actual SOS DP, O(n * 2^n) - use this one
for (int i = 0; i < (1 << n); i++)
    sos[i] = a[i];
for (int i = 0; i < n; i++)
    for (int mask = 0; mask < (1 << n); mask++)
        if (mask & (1 << i))
            sos[mask] += sos[mask ^ (1 << i)];
\end{lstlisting}
\subsection{Knapsack}
\begin{lstlisting}
// max value that fits in a knapsack of capacity W
int knapsackRec(int W, vector<int> &val, vector<int> &wt, int n,
                vector<vector<int>> &memo) {
    if (n == 0 || W == 0)
        return 0;
    if (memo[n][W] != -1)
        return memo[n][W];
    int pick = 0;
    // pick the nth item if it fits
    if (wt[n - 1] <= W)
        pick = val[n-1] + knapsackRec(W - wt[n-1], val, wt, n-1, memo);
    int notPick = knapsackRec(W, val, wt, n - 1, memo);
    return memo[n][W] = max(pick, notPick);
}
int knapsack(int W, vector<int> &val, vector<int> &wt) {
    int n = val.size();
    vector<vector<int>> memo(n + 1, vector<int>(W + 1, -1));
    return knapsackRec(W, val, wt, n, memo);
}
\end{lstlisting}
\section{Graph Theory}
\subsection{SPFA - Optimal BF O(V * E)}
\begin{lstlisting}
int q[3001]= {0};// queue for node
// d[i] = shortest path from start to node i
int d[1001]= {0};
bool f[1001]= {0};
int a[1001][1001]= {0}; // adjacency list
int w[1001][1001]= {0}; // adjacency matrix
int main(void) {
  int n=0, m=0;
  cin >> n >> m;
  for (int i=1; i<=m; i++){
  int x=0, y=0, z=0;
cin >> x >> y >> z;
// node x to node y has weight z
a[x][0]++;
 a[x][a[x][0]]=y;
  w[x][y]=z;
/*
// for undirected graph
 a[x][0]++;
 a[y][a[y][0]]=x;
w[y][x]=z;
*/
}
int s=0, e=0;
cin >> s >> e; // s: start, e: end
SPFA(s);
cout << d[e] << endl;
return 0;
}
void SPFA(int v0){
    int t,h,u,v;
    for (int i=0; i<1001; i++) d[i]=INT_MAX;
    for (int i=0; i<1001; i++) f[i]=false;
    d[v0]=0;
    h=0;
    t=1;
    q[1]=v0;
    f[v0]=true;
    while (h!=t){
        h++;
        if (h>3000) h=1;
        u=q[h];
        for (int j=1; j<=a[u][0]; j++){
            v=a[u][j];
            if (d[u]+w[u][v]<d[v]) // change to > if calculating longest path
            {
                d[v]=d[u]+w[u][v];
                if (!f[v]){
                    t++;
                    if (t>3000) t=1;
                    q[t]=v;
                    f[v]=true;
                }
            }
        }
        f[u]=false;
    }
}
\end{lstlisting}
\subsection{Dijkstra O(V + ElogV)}
\begin{lstlisting}
typedef pair<int, int> pairi;
int N = 20000 + 5;
vector<vector<pairi>> adj(N);
vector<int> dis(N, inf), parent(N);

void dijkstra(int src) {
    priority_queue<pairi, vector<pairi>, greater<pairi>> pq;
    dis[src] = 0;
    pq.push({0, src});
    while (pq.size()) {
        auto top = pq.top();
        pq.pop();
        for (auto i : adj[top.second]) {
            int v = i.first;
            int wt = i.second;
            if (dis[v]>dis[top.second]+wt) {
                dis[v]=dis[top.second]+wt;
                pq.push({dis[v], v});
            parent[v] = top.second;
            }
        }
    }
}
\end{lstlisting}
\subsection{Bellman-Ford O(V*E)}
\begin{lstlisting}
// needs: vector<pair<int,int>> adj[N]; int dist[N], parent[N];
// init dist[] = inf before calling, inf = 1e9
void bellmanFord(int num_of_nd, int src) {
    dist[src] = 0;
    for (int step = 0; step < num_of_nd; step++) {
        for (int i = 1; i <= num_of_nd; i++) {
            for (auto it : adj[i]) {
                int u = i;
                int v = it.first;
                int wt = it.second;
                if (dist[u] != inf && ((dist[u] + wt) < dist[v])) {
                    if (step == num_of_nd - 1) {
                        cout << "Negative cycle found\n";
                        return;
                    }
                    dist[v] = dist[u] + wt;
                    parent[v] = u;
                }
            }
        }
    }
    for (int i = 1; i <= num_of_nd; i++) cout << dist[i] << " ";
    cout << endl;
}
\end{lstlisting}
\subsection{Floyd-Warshall algorithm O(n\^{}3)}
\begin{lstlisting}
typedef double T;
typedef vector<T> VT;
typedef vector<VT> VVT;

typedef vector<int> VI;
typedef vector<VI> VVI;

bool FloydWarshall (VVT &w, VVI &prev){
  int n = w.size();
  prev = VVI (n, VI(n, -1));

  for (int k = 0; k < n; k++){
    for (int i = 0; i < n; i++){
      for (int j = 0; j < n; j++){
        if (w[i][j] > w[i][k] + w[k][j]){
          w[i][j] = w[i][k] + w[k][j];
          prev[i][j] = k;
        }
      }
    }
  }

  // check for negative weight cycles
  for(int i=0;i<n;i++)
    if (w[i][i] < 0) return false;
  return true;
}
\end{lstlisting}
\subsection{Topological sort}
\begin{lstlisting}
map<string, vector<string>> adj;
map<string, int> degree;
set<string> nodes;
vector<string> ans;
// adj: graph input, degree: cnt indegree,
// node: unique nodes, ans: path
int c = 0;
void topo_sort() {
    queue<string> qu;
// traverse all the nodes and check if its degree is 0 or not..
    for (string i : nodes) {
        if (degree[i] == 0) {
            qu.push(i);
        }
    }
    while (!qu.empty()) {
        string top = qu.front();
        qu.pop();
        ans.push_back(top);
        for (string i : adj[top]) {
            degree[i]--;
            if (degree[i] == 0) {
                qu.push(i);
            }
        }
    }
}
\end{lstlisting}
\subsection{Kruskal O(ElogE)}
\begin{lstlisting}
typedef pair<int, int> edge;

class Graph {
    vector<pair<int, edge>> G, T;
    vector<int> parent;
    int cost = 0;

   public:
    Graph(int n) {
        for (int i = 0; i < n; i++)
            parent.push_back(i);
    }

    void add_edges(int u, int v, int wt) { G.push_back({wt, {u, v}}); }

    int find_set(int n) {
        if (n == parent[n])
            return n;
        else
            return find_set(parent[n]);
    }

    void union_set(int u, int v) { parent[u] = parent[v]; }

    void kruskal() {
        sort(G.begin(), G.end());
        for (auto it : G) {
         int uRep=find_set(it.second.first);
        int vRep=find_set(it.second.second);
            if (uRep != vRep) {
                cost += it.first;
                T.push_back(it);
                union_set(uRep, vRep);
            }
        }
    }

    int get_cost() { return cost; }
    void print() {
        for (auto it : T)
            cout << it.second.first << " " << it.second.second
                 << "->" << it.first << endl;
    }
};

// g.add_edges(u, v, wt);
// g.kruskal();
\end{lstlisting}
\subsection{Prim - MST O(ElogV)}
\begin{lstlisting}
typedef pair<int, int> pii;

class Prims {
    map<int, vector<pii>> graph;
    map<int, int> visited;

   public:
    void addEdge(int u, int v, int w) {
        graph[u].push_back({v, w});
        graph[v].push_back({u, w});
    }

    vector<int> path(pii start) {
       vector<int> ans;
       // min-heap of (cost, node)
       priority_queue<pii, vector<pii>, greater<pii>> pq;
       pq.push({start.second, start.first});
       while (!pq.empty()) {
            pair<int, int> curr = pq.top();
            pq.pop();
            if (visited[curr.second])
                continue;
            visited[curr.second] = 1;
            ans.push_back(curr.second);
            for (auto i:graph[curr.second]){
               if (visited[i.first])
                   continue;
               pq.push({i.second, i.first});
            }
       }
       return ans;
    }
};
\end{lstlisting}
\subsection{Euler Tour - subtree flatten O(V+E)}
\begin{lstlisting}
unordered_map<int, int> Start, End, Val;
unordered_map<int, pair<int, int>> Range;
int start = 0;
void dfs(int node){
    visited[node] = true;
    Start[node] = start++;
    for (auto child : adj[node]){
        if (!visited[child])
            dfs(child);
    }
    End[node] = start - 1;
}
dfs(1);
vector<int> FlatArray(start + 5);
for (auto i : Start){
    FlatArray[i.second] = Val[i.first];
    Range[i.first] = {i.second, End[i.first]};
}
\end{lstlisting}
\subsection{LCA}
\begin{lstlisting}
// query O(logn)
// preprocessing O(n)
struct LCA {
    vector<ll> hi, euler, first, st;
    vector<bool> visit;
    ll n;
    LCA(vector<vector<ll>> &adj, ll root=0){
        n = adj.size();
        hi.resize(n);
        first.resize(n);
        euler.reserve(n * 2);
        visit.assign(n, false);
        dfs(adj, root);
        ll m = euler.size();
        st.resize(m * 4);
        build(1, 0, m - 1);
    }
    void dfs(vector<vector<ll>> &adj, ll node, ll h = 0) {
        visit[node] = true;
        hi[node] = h;
        first[node] = euler.size();
        euler.push_back(node);
        for (auto to : adj[node]) {
            if (!visit[to]) {
                dfs(adj, to, h + 1);
                euler.push_back(node);
            }
        }
    }
    void build(ll node, ll b, ll e) {
        if (b == e) {
            st[node] = euler[b];
        }
        else {
            ll mid = (b + e) / 2;
            build(node << 1, b, mid);
            build(node << 1 | 1, mid+1, e);
            ll l = st[node << 1], r = st[node << 1 | 1];
            st[node] = (hi[l]<hi[r])?l : r;
        }
    }
    ll query(ll node,ll b,ll e,ll L,ll R) {
        if (b > R || e < L)
            return -1;
        if (b >= L && e <= R)
            return st[node];
        ll mid = (b + e) >> 1;
        ll left=query(node << 1,b,mid,L,R);
        ll right = query(node << 1 | 1, mid + 1, e, L, R);
        if (left == -1)
            return right;
        if (right == -1)
            return left;
       return hi[left]<hi[right]?left:right;
    }

    ll lca(ll u, ll v) {
        ll left = first[u], right = first[v];
        if (left > right)
            swap(left, right);
        return query(1, 0, euler.size() - 1, left, right);
    }
};
\end{lstlisting}
\subsection{Min cost max flow}
\begin{lstlisting}
struct Edge{
    int from, to, capacity, cost;
};
vector<vector<int>> adj, cost, capacity;
const int INF = 1e9;
void shortest_paths(int n, int v0, vector<int>& d, vector<int>& p) {
    d.assign(n, INF);
    d[v0] = 0;
    vector<bool> inq(n, false);
    queue<int> q;
    q.push(v0);
    p.assign(n, -1);
    while (!q.empty()) {
        int u = q.front();
        q.pop();
        inq[u] = false;
        for (int v : adj[u]) {
            if (capacity[u][v] > 0 && d[v] > d[u] + cost[u][v]) {
                d[v] = d[u] + cost[u][v];
                p[v] = u;
                if (!inq[v]) {
                    inq[v] = true;
                    q.push(v);
                }
            }
        }
    }
}
int min_cost_flow(int N, vector<Edge> edges, int K, int s, int t) {
    adj.assign(N, vector<int>());
    cost.assign(N, vector<int>(N, 0));
    capacity.assign(N, vector<int>(N, 0));
    for (Edge e : edges) {
        adj[e.from].push_back(e.to);
        adj[e.to].push_back(e.from);
        cost[e.from][e.to] = e.cost;
        cost[e.to][e.from] = -e.cost;
        capacity[e.from][e.to] = e.capacity;
    }
    int flow = 0;
    int cost = 0;
    vector<int> d, p;
    while (flow < K) {
        shortest_paths(N, s, d, p);
        if (d[t] == INF)
            break;
        // find max flow on that path
        int f = K - flow;
        int cur = t;
        while (cur != s) {
            f = min(f, capacity[p[cur]][cur]);
            cur = p[cur];
        }
        // apply flow
        flow += f;
        cost += f * d[t];
        cur = t;
        while (cur != s) {
            capacity[p[cur]][cur] -= f;
            capacity[cur][p[cur]] += f;
            cur = p[cur];
        }
    }

    if (flow < K)
        return -1;
    else
        return cost;
}
\end{lstlisting}
\subsection{Biparitite}
\begin{lstlisting}
const int N=1000;
int adj[N][N];
int n,e;
bool isBicolored(int s){
    int colorArray[n];
    for(int i=0;i<n;i++)
        colorArray[i]=-1; //init no color;
    queue<int>q;
    q.push(s);
    colorArray[s]=1; //assigning first color
    while(!q.empty()){
        int senior = q.front();
        q.pop();
        if(adj[senior][senior]==1)
            return false;
        for(int i=0;i<n;i++){
            int junior=i;
            if(adj[senior][junior]==1){
                if(colorArray[junior]==colorArray[senior])   //successor(child/junior) having same color
                    return false;
                ///if(colorArray[junior]!=-1) continue;    ///not same color but have a color
                else if(colorArray[junior]==-1){        ///No color assigned
                    q.push(junior);
                    colorArray[junior]=!colorArray[senior]; ///assigning diff color
 }}}} return true;}
\end{lstlisting}
\section{Random Stuff}
\subsection{Knight Moves}
\begin{lstlisting}
int X[8]={2,1,-1,-2,-2,-1,1,2};
int Y[8]={1,2,2,1,-1,-2,-2,-1};
\end{lstlisting}
\subsection{Rand function}
\begin{lstlisting}
#define accuracy chrono::steady_clock::now().time_since_epoch().count()

mt19937 rng(accuracy);
int rand(int l, int r) {
   uniform_int_distribution<int> ludo(l, r);
   return ludo(rng);
}
\end{lstlisting}
\subsection{bit count in O(1)}
\begin{lstlisting}
int BitCount(unsigned int u){
     unsigned int uCount;
     uCount = u - ((u >> 1) & 033333333333) - ((u >> 2) & 011111111111);
     return ((uCount + (uCount >> 3)) & 030707070707) % 63;
}
\end{lstlisting}
\subsection{Matrix Exponentiation}
\begin{lstlisting}
#define vvi vector<vector<ll>>
ll n, m;
vvi matixMulti(vvi &a, vvi &b) {
    vvi res(n, vector<ll>(n, 0));
    for (ll i = 0; i < n; i++) {
        for (ll j = 0; j < n; j++) {
            for (ll k = 0; k < n; k++) {
                res[i][j] = (res[i][j] + (a[i][k] * b[k][j]) % mod) % mod;
            }
        }
    }
    return res;
}

vvi martixExp(vvi &base, ll power) {
    vvi identity(n, vector<ll>(n, 0));
    for (ll i = 0; i < n; i++)
        identity[i][i] = 1;

    while (power > 0) {
        if (power % 2) {
            identity = matixMulti(base, identity);
        }
        base = matixMulti(base, base);
        power /= 2;
    }
    return identity;
}
\end{lstlisting}
\subsection{sqrt decomposition(MO's Algo)}
\begin{lstlisting}
// https://www.spoj.com/problems/DQUERY/
#include <bits/stdc++.h>
using namespace std;
const int SIZE_1 = 1e6 + 10, SIZE_2 = 3e4 + 10;
class query{
public:
    int l, r, indx;
};
int block_size, cnt = 0;
int frequency[SIZE_1], a[SIZE_2];
void add(int indx){
    ++frequency[a[indx]];
    if (frequency[a[indx]] == 1)
        ++cnt;
}
void sub(int indx){
    --frequency[a[indx]];
    if (frequency[a[indx]] == 0)
        --cnt;
}
bool comp(query a, query b){
    if (a.l / block_size == b.l / block_size)
        return a.r < b.r;
    return a.l / block_size < b.l / block_size;
}
signed main(){
    int n; cin >> n;
    for(int i = 0; i < n; ++i) cin>>a[i];
    int q; cin >> q;
    int ans[q] = {};
    query Qur[q];
    for (int i = 0; i < q; ++i){
        int l, r; cin>>l>>r;

        Qur[i].l = l - 1;
        Qur[i].r = r - 1;
        Qur[i].indx = i;
    }
    block_size = sqrt(n); // sqrt(q) dileo hobe, but n is more accurate
    sort(Qur, Qur + q, comp);

    int ML = 0, MR = -1;
    for(int i = 0; i < q; ++i) {
        int L = Qur[i].l;
        int R = Qur[i].r;
        // fixing right pointer
        while (MR < R) add(++MR);
        while (MR > R) sub(MR--);
        // fixing left pointer
        while (ML < L) sub(ML++);
        while (ML > L) add(--ML);
        ans[Qur[i].indx] = cnt;
    }
    for (int i = 0; i < q; ++i)
        cout << ans[i] << '\n';
}//sqrt(n)
\end{lstlisting}
